\part{Number of types}\label{part:types}
In this part we prove the following result.
As we believe it to be of independent interest also in the context of infinite graphs,
all considered graphs may be finite or infinite in \Cref{part:types},
and we track the necessary finiteness assumptions in the arguments below.
Everywhere else in the paper, graphs are assumed to be finite.

% lax begin Lax185574.TypeBound
{\typeboundgeneraltrue\thmtypes*}
% lax end


We prove a stronger statement, which replaces the assumption of monadic dependence by the following  notion, which is a priori weaker.

% lax begin Lax185574.DpMinimality
Fix two formulas $\phi(\bar x,z),\psi(\bar y,z)$ and a graph $G$.
A \emph{clean $(\phi,\psi)$-pattern} of order $s\times t$ in $G$ consists of tuples 
$\bar a_1,\ldots,\bar a_s\in V(G)^{|\bar x|}$, $\bar b_1,\ldots,\bar b_t\in V(G)^{|\bar y|}$
and vertices $c_{i,j}$ for $i\in[s]$ and $j\in[t]$,
such that the following hold for all $i,i'\in[s],j,j'\in[t]$:
\begin{enumerate}
    \item $G\models \phi(\bar a_{i'},c_{i,j})$ if and only if $i=i'$,
    \item $G\models \psi(\bar b_{j'},c_{i,j})$ if and only if $j=j'$.
\end{enumerate}

\begin{definition}
We say that a class $\CC$ of structures is \emph{dp-minimal}
if for every two formulas $\phi(\bar x,z),\psi(\bar y,z)$ there is some $t\in\N$ 
such that no structure in $\CC$ contains a clean $(\phi,\psi)$-pattern of order $t\times t$.
\end{definition}
% lax end


% lax begin Lax185574.DpMinimalTypeBound
\begin{restatable}{theorem}{dpmintypes}\label{thm:type-bound-dp}
Let $\CC$ be a dp-minimal graph class and $k,q,c\in\N$.
For every \iftypeboundgeneral(possibly infinite) \fi graph $G \in \CC$, expansion $G^+$ of $G$ by $c$ unary predicates, and \iftypeboundgeneral finite \fi subset $U \subseteq V(G)$,
\begin{align}\label{eq:type-bound-dp}
    \left|\Types^{G^+}_{k,q}(V(G)/U)\right| \leq O_{\CC,q,k,c}(1)\cdot \poly_{\CC,q,k}(|U|).
\end{align}
    

\end{restatable}
% lax end

It is well-known and easy to see that monadic dependence implies dp-minimality (see \Cref{lem:mondep-dpmin}). Thus, \Cref{thm:type-bound-dp} implies \Cref{thm:type-bound}.
    For general classes of structures, however, dp-minimality is a strictly weaker assumption than monadic dependence\footnote{For example, $(\mathbb Z,+)$ is dp-minimal \cite[paragraph before Proposition 6.6]{ADHMS2016} but not monadically dependent:
    For $B=\setof{2^i}{i\in\mathbb N}$, the function $\set{u,v}\mapsto u+v$ defines an injection $f\from {B\choose 2} \to \mathbb Z$.
    Let $E\subset {B\choose 2}$ be arbitrary.
    In the unary expansion of $(\mathbb Z,+)$ with unary predicate $U$ marking $f(E)$,
    the formula $\phi(x,y):=U(x+y)$, restricted to $B$, defines the graph $(B,E)$. As $E$ is arbitrary and $\phi$ is fixed, $(\mathbb Z,+)$ is not monadically dependent.
    }. \Cref{thm:type-bound-dp} implies that in the setting of graphs, the two notions coincide.
% lax begin Lax185574.DpMinimalEquivalence
{\typeboundgeneraltrue    
\begin{corollary}\label{cor:dp-equivalence}
    Let $\CC$ be a class of \iftypeboundgeneral(possibly infinite) graphs.
    The following conditions are equivalent.
    \begin{enumerate}
        \item\label{it:mon-dep} $\CC$ is monadically dependent,
        \item\label{it:dp-min} $\CC$ is dp-minimal.
%         \item\label{it:type-bound} 
% For every $k,q,c\in\N$ and \fi graph $G \in \CC$, expansion $G^+$ of $G$ by $c$ unary predicates, and \iftypeboundgeneral finite \fi subset $U \subseteq V(G)$,
% \[
%     \left|\Types^{G^+}_{k,q}(V(G)/U)\right| \leq O_{\CC,q,k,c}(1)\cdot \poly_{\CC,q,k}(|U|).
% \]
    \end{enumerate}
\end{corollary}
}
% lax end
We first prove the corollary.
% lax begin Lax185574Proofs.DpMinimalEquivalence.monadicallyDependent_iff_dpMinimal
\begin{proof}
    \eqref{it:mon-dep}$\rightarrow$\eqref{it:dp-min} is an easy and well-known fact, proven in \Cref{lem:mondep-dpmin}.
    We prove the converse implication, 
    \eqref{it:dp-min}$\rightarrow$\eqref{it:mon-dep}.

    Let $\CC$ be a dp-minimal graph class; then $\CC$ satisfies the bound \eqref{eq:type-bound-dp} by \Cref{thm:type-bound-dp}.
    We argue that this bound implies monadic dependence.
    Indeed, let $c\in\N$, and let $\phi(x,y)$ be any formula in the signature of $c$-colored graphs, and let $q$ denote its quantifier rank. Let $N$ be a number and $H$ be the bipartite graph with parts $A=\{a_1,\ldots,a_N\}$ and $B=\setof{b_I}{I\subset [N]}$, where $a_i$ is adjacent to $b_I$ if and only if $i\in I$. Suppose there is some $G\in\CC$ and $c$-coloring $G^+$ of $G$, such that $H$ is an induced subgraph of $\phi(G^+)$, the graph with vertex set $V(G)$ and edge set $\setof{uv\in \binom{V(G)}{2}}{G^+\models \phi(u,v)\lor \phi(v,u)}$. Then $\tp^q(b_I/A)$ is different for every $I\subset [N]$,
so the number of types of quantifier rank $q$ over $A$ is at least $2^N$.
On the other hand, the bound in \Cref{thm:type-bound-dp} is at most polynomial in $|A|=N$, which implies that $N$ is bounded by a constant depending on $\CC$ and $\phi$.
Hence, not arbitrary graphs can be encoded in~$\CC$ by~$\phi$. 
\end{proof}
% lax end
\Cref{cor:dp-equivalence} extends to classes of structures over a finite\sz{necessary?} signature consisting of binary relations, via a standard encoding of such structures as graphs (see \cite[full version, Section 19]{flip-breakability}).

% lax begin Lax185574Proofs.DpMinimalTypeBound.typeCountOver_le
The rest of this part proves \Cref{thm:type-bound-dp}.
We state the intermediate results for monadically dependent classes, the setting of our main theorem.
Their proofs use this assumption only through \Cref{lem:phi-psi-patterns} and through the fact that the $s$-definable flips of graphs from $\CC$ have bounded VC-dimension.
Both hold for every dp-minimal class, by \Cref{lem:dpmin-phipsi-patterns,obs:dpmin-vc,cor:dpmin-flips}.
Hence, the same proof gives \Cref{thm:type-bound-dp}.
% lax end

% We make two remarks, which are of independent interest in the context of model theory.
% \begin{remark}\label{rem:converse}
% Conversely, the conclusion of \Cref{thm:type-bound} implies monadic dependence. 
% Indeed, let $c\in\N$, and let $\phi(x,y)$ be any formula in the signature of $c$-colored graphs, and let $q$ denote its quantifier rank. Let $N$ be a number and $H$ be the bipartite graph with parts $A=\{a_1,\ldots,a_N\}$ and $B=\setof{b_I}{I\subset [N]}$, where $a_i$ is adjacent to $b_I$ if and only if $i\in I$. Suppose there is some $G\in\CC$ and $c$-coloring $G^+$ of $G$, such that $H$ is an induced subgraph of $\phi(G^+)$, the graph with vertex set $V(G)$ and edge set $\setof{uv\in \binom{V(G)}{2}}{G^+\models \phi(u,v)\lor \phi(v,u)}$. Then $\tp^q(b_I/A)$ is different for every $I\subset [N]$,
% so the number of types of quantifier rank $q$ over $A$ is at least $2^N$.
% On the other hand, the bound in \Cref{thm:type-bound} is at most polynomial in $|A|=N$, which implies that $N$ is bounded by a constant depending on $\CC$ and $\phi$.
% Hence, not arbitrary graphs can be encoded in~$\CC$ by~$\phi$. 
% \end{remark}


% \begin{remark}
%      An inspection of our proof of \Cref{thm:type-bound} shows that it only uses the assumption that $\CC$ is a \emph{dp-minimal} class of graphs (see \Cref{app:patterns} for a definition).
%     Thus, the conclusion of \Cref{thm:type-bound} holds under the assumption of dp-minimality.

%     Monadic dependence implies dp-minimality, as shown in \Cref{lem:mondep-dpmin}.
%     For general classes of structures, however, dp-minimality is a strictly weaker assumption than monadic dependence\footnote{For example, $(\mathbb Z,+)$ is dp-minimal \cite[paragraph before Proposition 6.6]{ADHMS2016} but not monadically dependent:
%     For $B=\setof{2^i}{i\in\mathbb N}$, the function $\set{u,v}\mapsto u+v$ defines an injection $f\from {B\choose 2} \to \mathbb Z$.
%     Let $E\subset {B\choose 2}$ be arbitrary.
%     In the unary expansion of $(\mathbb Z,+)$ with unary predicate $U$ marking $f(E)$,
%     the formula $\phi(x,y):=\exists z. U(z)\land (x+y=z)$, restricted to $B$, defines the graph $(B,E)$. As $E$ is arbitrary and $\phi$ is fixed, $(\mathbb Z,+)$ is not monadically dependent.
%     }.
%      With the previous remark, we conclude that for classes $\CC$ of graphs, the notions of dp-minimality and monadic dependence coincide. This extends to classes of structures over a finite  signature consisting of binary relations, via the standard encoding of such structures as graphs (see \cite[full version, Section 19]{flip-breakability}).
% \end{remark}

\section{Closures}\label{sec:closures}
The goal of this section is to prove the following theorem, which can be seen as a finitary analogue of \cite[Theorem~7.1 in the full version]{flippers}.

% lax begin Lax185574.ClosureTheorem
\begin{restatable}[Closure Theorem]{theorem}{thmClosures}\label{thm:closures}
For every monadically dependent graph class $\CC$ and radius $r\in\N$ there are some $c,s\in\N$ with the following property.

\smallskip\noindent
For every $G \in \CC$ and finite set $A \subseteq V(G)$ there is a set $A^*\subseteq V(G)$ of size $O_{\CC,r}(|A|^c)$ such that for every vertex $v \notin A$ there is a set $S(v)\subseteq A^*$ of size at most $s$
and an $S(v)$-flip $H$ with
\[
    \dist^H(v,A) > r.
\]
\end{restatable}
% lax end



We also extend the statement to tuples.
For simplicity, and in contrast to \Cref{thm:closures}, the conclusion uses non-definable flips. 
Definability can easily be restored using the methods from~\cite{DBLP:conf/lics/BonnetDGKMST22} (see also~\cite{bonnet2025separability-icalp}).

% lax begin Lax185574.ClosureTheoremTuples
\begin{restatable}[Closure Theorem for Tuples]{theorem}{thmClosuresTuples}\label{thm:closuresTuples}
For every monadically dependent graph class $\CC$, radius $r\in\N$, and tuple size $k \in \N$, there are some $c,s\in\N$ with the following property.

\smallskip\noindent
For every $G \in \CC$ and finite set $A \subseteq V(G)$ there is a set $\HH$ of $O_{\CC,r,k}(|A|^c)$ many $s$-flips of $G$ such that for every tuple $\bar v \in (V(G) \setminus A)^k$ there is a flip $H \in \HH$ with
\[
    \dist^H(\bar v,A) > r.
\]
\end{restatable}
% lax end

\subsection{Preliminaries}


% lax begin Lax185574.DefinableFlips
\paragraph*{Definable flips.}
% lax begin Lax185574.AtomicTypes
For $s \in \N$, let $\Pi=\set{(x=y), (x\neq y)\land E(x,y),(x\neq y)\land \neg E(x,y)}$ denote the set of three possible atomic types of a pair of vertices in a graph.
For an $s$-tuple $\bar p=(p_1,\ldots,p_s)$ of vertices in a graph $G$ and a vertex $v$,
we identify $\atp(v/\bar p)$ with an $s$-tuple in $\Pi^s$, whose $i$th component is the atomic type of the pair $(v,p_i)$.
% lax end
For a symmetric relation $R \subseteq \Pi^s \times \Pi^s$, we denote by $G \oplus (\bar p, R)$ the $\set{p_1,\ldots,p_s}$-flip $H$ of $G$ obtained by partitioning $V(G)$ into atomic types over $\bar p$ and then flipping according to $R$.
More precisely, for $uv\in \binom{V}{2}$,
we have that $uv\in E(H)\triangle E(G)$ (where $\triangle$ denotes the symmetric difference) if and only if $(\atp^G(u/\bar p), \atp^G(v/ \bar p)) \in R$.
We call such a flip \emph{$s$-definable}.
Equivalently, an $s$-definable flip is an $S$-flip for some $S\subseteq V(G)$ of size at most $s$.
Note that an $s$-definable flip is also $s'$-definable for every $s'\ge s$.
% lax end


% lax begin Lax185574.DefinableFlipFormula
% lax begin Lax185574Proofs.DefinableFlipFormula.exists_isQF_flipOver_formula
\begin{observation}\label{obs:sim-def-flip}
    We can define the edges of a definable flip in the original graph.
    More precisely, for every $s\in\N$ and symmetric relation $R \subseteq \Pi^s \times \Pi^s$ there exists a quantifier-free formula $E_{\oplus R}(x,y,\bar z)$ such that for every graph $G$, tuple $\bar p \in V(G)^s$, and vertices $u,v\in V(G)$ we have 
    \[
        G\oplus(\bar p, R) \models E(u,v)
        \iff
        G \models E_{\oplus R}(u,v,\bar p).
    \]
\end{observation}
% lax end
% lax end

% lax begin Lax185574.FlipsMonadicallyDependent
% lax begin Lax185574Proofs.FlipsMonadicallyDependent.monadicallyDependent_definableFlipClass
The following is an easy consequence of the characterization of monadic dependence through transductions~(see, e.g., \cite{flip-breakability}). It also holds in the greater generality of $s$-flips (rather than $s$-definable flips).
% lax end
\begin{fact}\label{obs:flips-mon-dep}
    Let $\CC$ be a monadically dependent graph class and $s\in\N$.
    The class 
    consisting of all $s$-definable flips of graphs from $\CC$ is monadically dependent.
\end{fact}
% lax end


% lax begin Lax185574Proofs.DefinableFlipTransitivity.isDefinableFlip_trans
The following is well-known, see, e.g., Lemma~2.6 in the full version of~\cite{flippers}.
% lax end

% lax begin Lax185574.DefinableFlipTransitivity
\begin{fact}[Transitivity of definable flips]\label{lem:definable-flip-transitive}
    Let $G_1$ be an $s_1$-definable flip of $G_0$ and $G_2$ be an $s_2$-definable flip of $G_1$. Then $G_2$ is an $(s_1 + s_2)$-definable flip of $G_0$.
\end{fact}
% lax end


\subsection{The main argument}\label{sec:closures-main}

The argument is a finitary analogue of the reasoning in \cite[Section 7]{flippers}.
We use two lemmas. The first enlarges a given set $A$ to a set $A^*$ such that balls avoiding $A^*$ have few types over $A$. The second uses this bound to increase the separation radius by one.

% lax begin Lax185574.BallTypes
\begin{restatable}{lemma}{lemBallTypes}\label{lem:ball-types}
    For every monadically dependent graph class $\CC$ and $r,s\in\N$, there are some $k,t\in\N$ with the following property.
    For every $G\in\CC$ and finite set $A\subseteq V(G)$, there is $A^*\supseteq A$ of size $O_{\CC,r,s}(|A|^k)$ such that, for every $s$-definable flip $H$ of $G$ and every vertex $v$,
    \[
        \Ball_r^H(v)\cap A^*=\emptyset
        \quad\Longrightarrow\quad
        |\ATypes^H(\Ball_r^H(v)/A)|\leq t.
    \]
\end{restatable}
% lax end

The inequality above
can be equivalently stated as follows:
 $\Ball_r^H(v)$ intersects at most $t$ parts of 
the partition of $V(G)$ into atomic types over $A$ in $H$.
The proof of \Cref{lem:ball-types} is deferred to \Cref{sec:few-types-in-balls}.


% lax begin Lax185574.BallFlip
\begin{restatable}{lemma}{lemBallFlip}\label{lem:ball-flip}
For all $d,t\in \N$, there is some $z\in\N$ with the following property.

\smallskip\noindent
For every graph $G$ of VC-dimension at most $d$, radius $r\in\N$, finite set $A\subseteq V(G)$, and arbitrary set $C\subseteq V(G)$ satisfying
\begin{center}
    $\Ball_r^G(C) \cap A = \emptyset$
    \qquad and \qquad
    $|\ATypes^G(\Ball_r^G(C)/A)|\leq t$,
\end{center}
there exists a $z$-definable flip $H$ of $G$ in which $\Ball_{r+1}^H(C) \cap A = \emptyset$.
\end{restatable}
% lax end

The proof is deferred to \Cref{sec:definable-flips-for-balls}. We now conclude \Cref{thm:closures} from the two lemmas.

% lax begin Lax185574.ClosureTheorem
\thmClosures*
% lax end

For brevity, we say a flip $H$ of $G$ \emph{$r$-separates} two sets $A, B\subseteq V(G)$ if $\dist^H(A,B)>r$.
This terminology extends from sets to vertices as expected. 

% lax begin Lax185574Proofs.ClosureTheorem.closure_theorem
\begin{proof}[Proof of \Cref{thm:closures}]
    We induct on $r$, with $A^*\coloneqq A$, $S(v)\coloneqq \emptyset$, and $H\coloneqq G$ proving the base case $r=0$. 
    Suppose the theorem holds at radius $r$, with parameter bound $s$. 
    By \Cref{obs:flips-mon-dep}, and since monadically dependent classes have bounded VC-dimension (\Cref{obs:md-vc}), it follows that
    every $s$-definable flip of a graph in $\CC$ has VC-dimension at most some $d\leq O_{\CC,r}(1)$.

    Fix $G\in\CC$ and a finite set $A\subseteq V(G)$. Assume $A$ is nonempty, as otherwise the statement is trivial.
    Applying \Cref{lem:ball-types} to $r,s,G,A$ gives a polynomial-size superset $A^*\supseteq A$ and a constant $t \leq O_{\CC,r}(1)$ such that every $r$-ball in any $s$-definable flip $H$ of $G$ that avoids $A^*$ realizes at most $t$ types over $A$ in $H$.
    (Crucially, to avoid cyclic dependencies, note that $s$ is an input, not an output, of \Cref{lem:ball-types}.)

    Apply the induction hypothesis to $A^*$, obtaining a vertex set $B$ of size $\poly_{\CC,r}(|A^*|)\leq\poly_{\CC,r}(|A|)$ such that every vertex outside $A^*$ is $r$-separated from $A^*$ by a flip defined over a subset of $B$ of size at most $s$. We may assume that $B$ is nonempty.


    For each $\bar p\in B^s$ and each symmetric relation $R\subseteq\Pi^{s}\times\Pi^{s}$, put $H\coloneqq G\oplus(\bar p,R)$.
    Let $V_H$ be the set of vertices $v$ that are $r$-separated from $A^*$ by $H$, that is, $\Ball_r^H(v)\cap A^*=\emptyset$.

    We assign to each $v \in V_H$ its \emph{profile}
    \[
        \mathrm{profile}_H(v) \coloneqq \ATypes^H(\Ball_r^H(v)/A).
    \]
Equivalently,
if  $\PP_A$ denotes the partition of $V(G)$ into atomic types over $A$ in $H$,
then $\mathrm{profile}_H(v)=
\Ball_r^H(v)/\PP_A$
is the set of equivalence classes that are intersected by $\Ball_r^H(v)$.


Then $\mathrm{profile}_H(v)$ 
is a subset of $\ATypes^H(V(G)/ A)$ of size at most $t$, 
by \Cref{lem:ball-types}.
By the Sauer--Shelah lemma (\Cref{lem:ss}), we have that
\[
    |\ATypes^H(V(G)/ A)|\leq (|A|+1)^d+|A|.
\]
It follows that 

    \[
        |\{\mathrm{profile}_H(v) \,:\, v \in V_H\}|
        \leq
        |\ATypes^H(V(G)/ A)|^t \leq \big((|A|+1)^d+|A|\big)^t 
        \leq 
        \poly_{\CC,r}(|A|).
    \]
    
    Group the vertices of $V_H$ by their profiles.
    For each group $C$, the ball $\Ball_r^H(C)$:
    \begin{itemize}
        \item avoids $A$, because each $\Ball_r^H(v)$ with $v \in V_H$ does, and
        \item realizes at most $t$ types over $A$ in $H$, because all $v \in C$ have the same profile of size ${\leq}t$.    
    \end{itemize}
    Applying \Cref{lem:ball-flip} yields a 
    number $z\le O_{d,t}(1)\le O_{\CC,r}(1)$ and a
    further $z$-definable flip $H_C$ of $H$ that $(r+1)$-separates $C$ from $A$.
    By transitivity of definable flips (\Cref{lem:definable-flip-transitive}), $H_C$ is an $(s+z)$-definable flip of $G$.
    Choose a set $U_{\bar p,R,C}\subseteq V(G)$ of size at most $s+z$ such that $H_C$ is a $U_{\bar p,R,C}$-flip of $G$.

    We finish the proof by setting
    \[
        A^{**}\coloneqq A^*\cup\bigcup_{\bar p,R,C}U_{\bar p,R,C}.
    \]
    There are at most $|B|^s$ choices of $\bar p$, $O_s(1)$ choices of $R$ for each $\bar p$, and $\poly_{\CC,r}(|A|)$ groups $C$ for each choice of $\bar p,R$.
    Since each $U_{\bar p,R,C}$ has size at most $s+z$, we have $|A^{**}|\leq\poly_{\CC,r}(|A|)$.

    For each $v\notin A^*$, the induction hypothesis gives a tuple $\bar p \in B^s$ and a $\bar p$-flip $H$ that $r$-separates $v$ from $A^*$.
    Let $R$ be its flip relation, so that $H=G\oplus (\bar p,R)$, and let $C$ be the group of $V_H$ containing $v$.
    Then we may take $S(v)\coloneqq U_{\bar p,R,C}\subseteq A^{**}$, since $H_C$ $(r+1)$-separates $v$ from $A$.
    For $v\in A^*\setminus A$, take $S(v)\coloneqq \{v\}\subseteq A^{**}$ and the flip that isolates $v$.
    Hence, the theorem holds at radius $r+1$ with parameter set $A^{**}$ and parameter bound $s+z$.
\end{proof}
% lax end







    
\subsection{Few types in balls}\label{sec:few-types-in-balls}
In this subsection we prove \Cref{lem:ball-types}.
We use the following Ramsey-type result for set systems due to Ding, Oporowski, Oxley, and Vertigan~\cite{DBLP:journals/jct/DingOOV96}.

% lax begin Lax185574.AtomicMatching
\begin{definition}
    For a graph $G$ and tuples $\bar a, \bar b \in V(G)^t$, we say $\bar a = a_1 \ldots a_t$ and $\bar b = b_1 \ldots b_t$ are \emph{atomically matched} 
    if there is a relation ${\sim} \in \{=,\neq,\geq\}$ such that $G \models E(a_i, b_j)$ if and only if $i \sim j$ for all $i,j \in [t]$.
    In particular, the bipartite graph semi-induced between the tuples $\bar a$ and $\bar b$ is a matching, a co-matching, or a half-graph.
\end{definition}
% lax end

% lax begin Lax185574.MatchedTuples
% lax begin Lax185574Proofs.MatchedTuples.exists_atomicallyMatched
\begin{theorem}[{\cite[Corollary~2.4]{DBLP:journals/jct/DingOOV96}, see also~\cite[Theorem~2]{gravier2004}}]\label{thm:matching}
    There exists a computable function $Q:\N\rightarrow\N$ such that for every $\ell\in\N$ and for every bipartite graph $G = (L,R,E)$ with
    $|\ATypes(L/R)| > Q(\ell)$, there are two tuples $\bar a \in L^\ell$ and $\bar b \in R^\ell$ that are atomically matched in $G$.
\end{theorem}
% lax end
% lax end

We use the above result to construct the following obstructions to monadic dependence.

% lax begin Lax185574.PhiPsiPatterns
\begin{restatable}{definition}{defphipsi}\label{def:phi-psi-patterns}
    Fix first-order formulas $\phi(\bar x,z)$ and $\psi(\bar y,z)$. A \emph{$(\phi,\psi)$-pattern} of order $s \times t$ in a graph~$G$ consists of tuples $\bar a_1, \ldots, \bar a_s \in V(G)^{|\bar x|}$, $\bar b_1, \ldots, \bar b_t \in V(G)^{|\bar y|}$, and vertices ${(c_{i,j}: i\in [s],j\in [t])}$ such that the following holds
    for all $i,i'\in [s]$ and $j,j'\in[t]$:
    \begin{enumerate}
        \item $G\models \phi(\bar a_i,c_{i,j})$ 
        and $G\not\models \phi(\bar a_{i'},c_{i,j})$ if $i'>i$,
        \item $G\models \psi(\bar b_j,c_{i,j})$ 
        and $G\not\models \psi(\bar b_{j'},c_{i,j})$ if $j'>j$.
    \end{enumerate}
\end{restatable}
% lax end
\noindent We call $\setof{c_{i,j}}{i\in[s],j\in[t]}$ the \emph{grid-points} of the pattern, and $s$ its \emph{height}. See \Cref{fig:phi-psi-pattern}.
\begin{figure}[htbp]
    \definecolor{patternphi}{RGB}{0,90,200}
    \definecolor{patternpsi}{RGB}{200,40,40}
    \centering
    \begin{tikzpicture}[
        x=1.05cm,y=0.90cm,font=\small,
        grid point/.style={circle,fill=black,inner sep=1.5pt},
        % Tuples are drawn as o=o: two grid-point-sized dots (radius
        % 1.5pt*sqrt(2)) joined by a bar.
        a tuple/.style={rectangle,inner sep=0pt,
            minimum width=10pt,minimum height=5pt,
            path picture={
                \fill ([xshift=2.12pt]path picture bounding box.west) circle (2.12pt);
                \fill ([xshift=-2.12pt]path picture bounding box.east) circle (2.12pt);
                \fill ([xshift=2.12pt,yshift=-2.12pt]path picture bounding box.west)
                    rectangle ([xshift=-2.12pt,yshift=2.12pt]path picture bounding box.east);
            }},
        b tuple/.style={rectangle,inner sep=0pt,
            minimum width=5pt,minimum height=10pt,
            path picture={
                \fill ([yshift=-2.12pt]path picture bounding box.north) circle (2.12pt);
                \fill ([yshift=2.12pt]path picture bounding box.south) circle (2.12pt);
                \fill ([xshift=-2.12pt,yshift=-2.12pt]path picture bounding box.north)
                    rectangle ([xshift=2.12pt,yshift=2.12pt]path picture bounding box.south);
            }},
        phi incidence/.style={draw=patternphi,line width=0.5pt},
        psi incidence/.style={draw=patternpsi,line width=0.5pt}
    ]
        \foreach \i in {1,...,5} {
            \node[a tuple] (a-\i) at (-0.2,-\i) {};
        }
        \foreach \i/\aindex in {1/1,5/s} {
            \node[anchor=east,text=patternphi,inner sep=1pt] at ([xshift=-1pt]a-\i.west) {$\bar a_{\aindex}$};
        }
        \node[anchor=east,text=patternphi,inner sep=1pt] at ([xshift=-1pt]a-3.west) {$\bar a_i$};
        \foreach \j in {1,...,5} {
            \node[b tuple] (b-\j) at (\j,0.15) {};
        }
        \foreach \j/\bindex in {1/1,5/t} {
            \node[text=patternpsi] at (\j,0.65) {$\bar b_{\bindex}$};
        }
        \node[text=patternpsi] at (3,0.65) {$\bar b_j$};
        \foreach \i in {1,...,5} {
            \foreach \j in {1,...,5} {
                \coordinate (c-\i-\j) at (\j,-\i);
            }
        }
        % Curved fans separate the individual required formula incidences.
        \foreach \i in {1,...,5} {
            \foreach \j in {1,...,5} {
                \pgfmathsetmacro{\incidencewidth}{0.5+0.7*(\i==3)*(\j==3)}
                \draw[phi incidence,line width=\incidencewidth pt] (a-\i.east)
                    .. controls (0.45,-\i+0.05*\j)
                    and (\j-0.25,-\i+0.10+0.09*\j) .. (c-\i-\j);
                \draw[psi incidence,line width=\incidencewidth pt] (b-\j.south)
                    .. controls (\j-0.12-0.10*\i,-0.35)
                    and (\j-0.12-0.10*\i,-\i+0.35) .. (c-\i-\j);
            }
        }
        % Four illustrative phi-incidences with earlier row tuples.
        \foreach \earlier/\i/\j in {2/3/1,3/5/2,4/5/4} {
            \pgfmathsetmacro{\incidencewidth}{0.5+0.7*(\i==3)*(\j==3)}
            \draw[phi incidence,dashed,line width=\incidencewidth pt] (a-\earlier.east)
                .. controls (0.55,-\earlier)
                and (\j-0.40,-\i+0.40) .. (c-\i-\j);
        }
        \draw[phi incidence,dashed,line width=1.2pt] (a-2.east)
            .. controls (0.55,-1.75) and (2.60,-2.60) .. (c-3-3);
        % Five illustrative psi-incidences with earlier column tuples.
        \foreach \earlier/\i/\j in {1/1/2,1/2/3,2/1/3,4/2/5} {
            \pgfmathsetmacro{\incidencewidth}{0.5+0.7*(\i==3)*(\j==3)}
            \draw[psi incidence,dashed,line width=\incidencewidth pt] (b-\earlier.south)
                .. controls (\earlier-0.45,-0.40)
                and (\j-0.35,-\i+0.35) .. (c-\i-\j);
        }
        \draw[psi incidence,dashed] (b-1.south)
            .. controls (0.35,-0.65) and (1.65,-2.90) .. (c-3-2);
        \foreach \i in {1,...,5} {
            \foreach \j in {1,...,5} {
                \node[grid point] at (c-\i-\j) {};
            }
        }
        \foreach \i/\j/\rowindex/\colindex in {1/1/1/1,1/5/1/t,5/1/s/1,5/5/s/t} {
            \node[anchor=north west,fill=white,inner sep=1pt,font=\scriptsize]
                at (\j+0.10,-\i-0.17) {$c_{\rowindex,\colindex}$};
        }
        \node[anchor=north west,fill=white,inner sep=1pt,font=\scriptsize]
            at (3.10,-3.17) {$c_{i,j}$};
        \node[grid point,inner sep=2.8pt] at (c-3-3) {};

        % Truth profiles at the marked grid-point; ? means unrestricted.
        \node[anchor=west] at (6.0,-0.3) {At a grid-point $c_{i,j}$:};
        \node[anchor=west,text=patternphi] at (6.0,-1.1)
            {$\phi(\bar x,c_{i,j})$};
        \node[anchor=west,text=patternpsi] at (6.0,-3.3)
            {$\psi(\bar y,c_{i,j})$};
        \foreach \k/\profilex/\p/\q in {1/7/?/?,2/8.35/?/?,3/9/1/1,4/9.65/0/0,5/11/0/0} {
            \node[draw,minimum size=5mm,text=patternphi] (phi-box-\k)
                at (\profilex,-2.1) {$\p$};
            \node[draw,minimum size=5mm,text=patternpsi] (psi-box-\k)
                at (\profilex,-4.3) {$\q$};
        }
        \foreach \k/\aindex/\bindex in {1/1/1,2/i-1/j-1,3/i/j,4/i+1/j+1,5/s/t} {
            \node[font=\scriptsize] (phi-label-\k)
                at ($(phi-box-\k.center)+(0,0.45)$) {$\bar a_{\aindex}$};
            \node[font=\scriptsize,anchor=south] (psi-label-\k)
                at ($(phi-label-\k.south)+(0,-2.2)$) {$\bar b_{\bindex}$};
        }
        \foreach \leftindex/\rightindex in {1/2,4/5} {
            \node[font=\scriptsize]
                at ($(phi-box-\leftindex.center)!0.5!(phi-box-\rightindex.center)$) {$\cdots$};
            \node[font=\scriptsize]
                at ($(psi-box-\leftindex.center)!0.5!(psi-box-\rightindex.center)$) {$\cdots$};
        }
    \end{tikzpicture}
    \caption{A $(\phi,\psi)$-pattern of order $s\times t$ for $s=t=5$.
        % Horizontal ellipses represent the tuples $\bar a_i$, and vertical ellipses represent the tuples $\bar b_j$.
        Edges indicate $\phi(\bar a_{i'},c_{i,j})$ or $\psi(\bar b_{j'},c_{i,j})$.
        Each grid-point $c_{i,j}$ is related via $\phi$ to the tuple $\bar a_i$ in its row and possibly (indicated by dashed edges) to some earlier row tuples, and analogously for $\psi$ and the column tuples.
        % For each $c_{i,j}$, the truth values at later indices $i'>i$ and $j'>j$
        % must be $0$, while those at earlier indices are unrestricted, denoted by~$?$.
        The profiles on the right illustrate these conditions at a grid-point $c_{i,j}$.}
    \label{fig:phi-psi-pattern}
\end{figure}



\medskip
The following lemma is a corollary of~\cite[Lemma~C.1]{ordered-tww}. We include a self-contained proof in \Cref{app:patterns} for completeness.

% lax begin Lax185574.ForbiddenPhiPsiPatterns
\begin{restatable}{lemma}{lemphipsi}\label{lem:phi-psi-patterns}
    For every monadically dependent graph class $\CC$ and formulas $\phi(\bar x,z),\psi(\bar y,z)$, there is some $t\in\N$ such that no graph $G\in \CC$ contains a $(\phi,\psi)$-pattern of order $t \times t$.
\end{restatable}
% lax end



% lax begin Lax185574.PhiPatterns
\begin{definition}
    For a first-order formula $\phi(\bar x,z)$, a \emph{$\phi$-pattern} is a $(\phi,\psi)$-pattern
    with $\psi(y,z)=E(y,z)$ or $\psi(y,z)=\neg E(y,z)$.
    We call those a
    \emph{$(\phi,E)$-pattern} 
    and a \emph{$(\phi,\neg E)$-pattern}, respectively.
    Such a pattern consists of tuples $\bar a_1,\ldots,\bar a_s\in V(G)^{|\bar x|}$,
    vertices $b_1,\ldots,b_t\in V(G)$ and $(c_{i,j}:i\in[s],j\in [t])$.
    We call $\bar b=(b_1,\ldots,b_t)$
    the \emph{base} of the pattern. 
\end{definition}
% lax end

% lax begin Lax185574.ForbiddenPhiPatterns
% lax begin Lax185574Proofs.ForbiddenPhiPatterns.forbidden_phi_patterns
\begin{corollary}\label{lem:forbidden-patterns}
        For every monadically dependent graph class $\CC$ and formula $\phi(\bar x,z)$, there is some $t\in \N$ such that no graph $G\in \CC$ contains a $\phi$-pattern of order $t \times t$.
\end{corollary}
% lax end
% lax end


We now prove \Cref{lem:ball-types}, which we repeat for convenience.

% lax begin Lax185574.BallTypes
\lemBallTypes*
% lax end
% lax begin Lax185574Proofs.BallTypes.ball_types
\begin{proof}
    For each symmetric relation $R\subseteq\Pi^s \times \Pi^s$, let $\phi_R(\bar z x,y)$ be a first-order formula expressing that $x$ and $y$ have distance at most $r$ in the flip $G\oplus(\bar z,R)$, where $\bar z$ is an $s$-tuple.
Its existence follows from \Cref{obs:sim-def-flip}.

Let $k\ge 1$ be the bound from 
     \Cref{lem:forbidden-patterns}, such that no $G\in\CC$ contains a $\phi_R$-pattern of order $k\times k$, for any $R$.
    Let $Q\from\N\to\N$ be as in \Cref{thm:matching}, and set
    \[
        t\coloneqq Q\big(k \cdot |\Pi^s|^2\big)=Q(k\cdot 9^s).
    \]

    Fix $G\in\CC$ and a finite set $A\subseteq V(G)$.
    For every $R$ as above and $N\in \set{E,\neg E}$, and every $\bar b\in A^k$, choose a maximum-height $(\phi_R,N)$-pattern with base $\bar b$ in $G$. Each of them is an $\ell\times k$ pattern for some $\ell<k$ (possibly $\ell=0$).
    
    Let $A^*$ consist of $A$ and all grid-points of the chosen patterns.
    There are $O_s(1)$ many choices of $R$ and $|A|^k$ many choices of $\bar b$, and each chosen pattern contributes at most $k^2$ vertices, so
    \[
        |A^*|\leq |A|+O_s(1)\cdot|A|^k\cdot k^2
        =O_{\CC,r,s}(|A|^k).
    \]

    Fix a symmetric relation $R\subseteq\Pi^s \times \Pi^s$, tuple $\bar p \in V(G)^s$, and vertex $v \in V(G)$. Let
    $H=G\oplus(\bar p,R)$ and suppose $B\coloneqq \Ball_r^H(v)$ satisfies $B\cap A^*=\emptyset$.
    We show that $B$ realizes at most $t$ atomic types over $A$ in $H$, thus proving the lemma.
    
    Suppose the contrary.
    By \Cref{thm:matching}, there are tuples $\bar u'$ in $B$ and $\bar w'$ in $A$, each of length $k \cdot |\Pi^s|^2$, that are atomically matched in $H$.

    Let $\PP$ be the partition of $V(G)$ into atomic types over $\bar p$; then $H$ is a $\PP$-flip of $G$ and $|\PP|\le |\Pi^s|$.
    By the pigeonhole principle, we can pick corresponding subsequences $\bar u$ of $\bar u'$ and $\bar w$ of $\bar w'$ of length $k$,
    such that $\bar u$ 
    is contained in a single part $X\in \PP$
    and $\bar w$ is contained in a single part $Y\in\PP$.
    Write $\bar u=(u_1,\ldots,u_k)$ and $\bar w=(w_1,\ldots,w_k)$.
    Since $H$ is a $\PP$-flip of $G$, adjacency between $X$ and $Y$ is either preserved or complemented.
    Together with the matching condition in $H$, this gives a choice of $N\in\set{E,\neg E}$ such that
    \[
        \begin{aligned}
            G&\models N(w_j,u_j)&\text{for every }j\in[k],\\
            G&\not\models N(w_{j'},u_j)&\text{for }1\le j<j'\le k.
        \end{aligned}
    \]
    
    Consider the maximum-height $(\phi_R,N)$-pattern chosen for $R$ and $N$ and for  base $\bar b\coloneqq\bar w=(b_1,\ldots,b_k)$.
    Write $\bar a_1,\ldots,\bar a_\ell$ for its witnessing tuples and $c_{i,j}$ for its grid-points, where $i\in[\ell]$ and $j\in[k]$ (possibly $\ell=0$).
    All these grid-points belong to $A^*$.
    Append a new row by setting
    \[
        \bar a_{\ell+1}\coloneqq\bar p v,
        \qquad
        c_{\ell+1,j}\coloneqq u_j\quad\text{for every }j\in[k].
    \]
    We verify that this yields a $(\phi_R,N)$-pattern with the same base $\bar b$ and height $\ell+1$, contradicting the maximality of $\ell$.        
    The column condition holds by the choice of $N$, $b_j=w_j$ and $c_{\ell+1,j}=u_j$:
    \[
        \begin{aligned}
            G&\models N(b_j,c_{\ell+1,j})&\text{for every }j\in[k],\\
            G&\not\models N(b_{j'},c_{\ell+1,j})&\text{for }1\le j<j'\le k.
        \end{aligned}
    \]    
    For every $j\in[k]$, we have $G\models\phi_R(\bar a_{\ell+1},c_{\ell+1,j})$, since $c_{\ell+1,j}\in B=\Ball_r^H(v)$.
    Moreover, for every $i\in[\ell]$, we have $G\not\models\phi_R(\bar a_{\ell+1},c_{i,j})$, since $c_{i,j}\in A^*$ and $B\cap A^*=\emptyset$. This proves the lemma.
\end{proof}
% lax end




\subsection{Definable flips for balls}\label{sec:definable-flips-for-balls}

% lax begin Lax185574Proofs.VCDuality.vc_duality
The following lemma is a corollary of the $(p,q)$-Theorem for bounded VC-dimension from~\cite{Matousek:p-q-theorem}.
See~\cite[Corollary~6.13]{simon_book} or~\cite[full version, Theorem~E.2]{flip-width} for self-contained proofs.
% lax end

For a graph $G$ and sets $S,X\subseteq V(G)$, we say that $S$ \emph{dominates} $X$ if every vertex of $X \setminus S$ has a neighbor in $S$, and that $S$ \emph{antidominates} $X$ if every vertex of $X \setminus S$ has a nonneighbor in $S$.

% lax begin Lax185574.VCDuality
\begin{lemma}[Bounded-VC duality]\label{lem:vc-duality}
    For every $d\in\N$ there is some $k\in\N$ with the following property.
    Let $G$ be a graph of VC-dimension at most $d$, and let $X,Y\subseteq V(G)$ be two (possibly overlapping) sets with at least one of $X,Y$ finite.
    Then there exists a set $\sample(X,Y)$ of size at most $k$ such that at least one of the following holds:
    \begin{enumerate}
        \item $\sample(X,Y)$ is contained in $Y$ and dominates $X$;
        \item $\sample(X,Y)$ is  contained in $X$ and antidominates $Y$.
    \end{enumerate}
\end{lemma}
% lax end
We remark that in the sources we have given, the above lemma is stated either for finite set systems or finite binary relations.
The translation to finite graphs is straightforward, by either viewing the graph as a set system of neighborhoods or working with the adjacency relation.
Moreover, the statement lifts to infinite graphs if at least one of the sets $X$ and $Y$ is finite, as the other one can be replaced by its finite subset representing all neighborhoods in the  finite one.
(However, the statement may fail if both $X$ and $Y$ are infinite, as witnessed by the bi-infinite half-graph with vertices $\mathbb{Z} \times \{1,2\}$ and edges between $(x,1)$ and $(y,2)$ if $x\leq y$.)

It was observed already in~\cite{DBLP:conf/lics/BonnetDGKMST22} that \Cref{lem:vc-duality} is useful to obtain definable flips.
We will use it in a similar way to deduce the remaining result of \Cref{sec:closures-main}.

% lax begin Lax185574.BallFlip
\lemBallFlip*
% lax end




Note that, given the assumptions of the lemma, it is easy to find a (non-definable) $(2^t+t+1)$-flip $H'$ of $G$ in which $\Ball_{r+1}^{H'}(C) \cap A = \emptyset$: group the vertices of $A$ and of $\Ball_r^G(C)$ by their neighborhoods in the other set, which gives at most $2^t$ and $t$ groups, respectively, and flip every group of $A$ with every group of $\Ball_r^G(C)$ that it is complete to.
So the challenge is really to get such a flip in a definable way.

We also remark that \Cref{lem:ball-flip} can be seen as a finitization of Lemma~7.3 in the full version of~\cite{flippers}.
That lemma works in a more model-theoretic setting, where $A$ is a (possibly infinite) model instead of a set.
It uses a version of \Cref{lem:vc-duality} which also works for infinite $X,Y$, under the stronger 
assumption that $G$ excludes a half-graph of order $d$.

% lax begin Lax185574Proofs.BallFlip.exists_ball_flip
\begin{proof}[Proof of \Cref{lem:ball-flip}]
    Put $B\coloneqq \Ball_r^G(C)$.
    Let $\AA$ be the partition of $A$ that groups vertices by their neighborhood in $B$, and let $\BB$ be the partition of $B$ that groups vertices by their neighborhood in $A$.
    By assumption $|\BB|\leq t$.
    The neighborhood in $B$ of a vertex of $A$ is a union of parts of $\BB$, so $|\AA|\leq2^t$.
    Moreover, every pair $P\in\AA$, $Q\in\BB$ is homogeneously connected: is either complete or anticomplete.

    Let $k \leq O_d(1)$ be given by \Cref{lem:vc-duality}.
    Pick any set $S$ which contains
    \begin{itemize}
        \item an arbitrary representative $\rep(P)\in P$ for each part $P\in\AA\cup\BB$, and
        \item
        the set $\sample(X,Y)$ supplied by \Cref{lem:vc-duality}, for every $X\in\AA$ and $Y\in\AA\cup\{V(G)\setminus B\}$         
        (the lemma is applicable because $X \subseteq A$ is finite).
    \end{itemize}
    We can pick such a set with $|S|\le z$ for some $z\le O_{k,t}(1)\le O_{d,t}(1)$.
    We will construct $H$ as an $S$-flip of $G$; in particular, $H$ is $z$-definable.
    
    Let $\PP_S$ be the partition of $V(G)$ into atomic types over $S$.
    Write $x\sim y$ when $x$ and $y$ are in the same $\PP_S$-part.
    Observe the following two properties of $\PP_S$.

    \begin{claim}
    $\PP_S$ refines $\AA$ on $A$ and $\BB$ on $B$.
    \end{claim}
    \begin{claimproof}
        If two vertices $a$ and $a'$ of $A$ are in different $\AA$-parts, then their adjacencies differ on some vertex $b \in B$.
        Let $P$ be the $\BB$-part containing $b$.
        The vertex $\rep(P) \in S$ has the same adjacencies to $A$ as $b$ and hence distinguishes $a$ and $a'$ in $\PP_S$.
        Thus, $\PP_S$ refines $\AA$ on $A$ and symmetrically also $\BB$ on $B$.
    \end{claimproof}

    \begin{claim}\label{clm:dagger}
    Let $X\in \AA$ and $Y\in \AA\cup\set{V(G)-B}$. 
    There do not exist vertices $x,y,x',y'$ with:
    \begin{equation*}\tag{$\dagger$}\label{eq:ball-flip-crossing}
        x\in X,\quad y\in Y,\quad x\sim x',\quad y\sim y',
        \quad Y\subseteq N^{\overline G}(x'),\quad X\subseteq N^G(y').
    \end{equation*}    
    \end{claim}
    \begin{claimproof}
    Suppose that $\sample(X,Y)$ lies in $Y$ and dominates $X$. 
    Then there is some  $s\in\sample(X,Y)\cap Y$ with $sx\in E(G)$.
    If
    $Y\subseteq N^{\overline G}(x')$,
then $sx'\notin E(G)$.
Thus, $x\nsim x'$.

The case when $\sample(X,Y)$ lies in $X$ and antidominates $Y$, is symmetric, by swapping $X$ with $Y$ and $G$ with $\overline G$.
    Both cases are illustrated in \Cref{fig:ball-flip-crossing}.
    \end{claimproof}
    \begin{center}
        \captionsetup{type=figure}
        \definecolor{bfAnti}{RGB}{200,40,40}%
\definecolor{bfHl}{RGB}{0,90,200}%
{\setlength{\fboxsep}{6pt}\makebox[\linewidth]{\colorbox{white}{%
\begin{tikzpicture}[
        vertex/.style={circle, fill=black, inner sep=0pt, minimum size=3.2pt},
        blob/.style={draw=black!70, line width=0.6pt, fill=black, fill opacity=0.04},
        complete/.style={draw=black, line width=0.6pt},
        anticomplete/.style={draw=bfAnti, dash pattern=on 2.5pt off 2pt, line width=0.6pt,
            shorten <=2pt},
        sameType/.style={black!60, densely dotted, line width=0.7pt},
        lbl/.style={inner sep=1pt},
        sampleSet/.style={draw=bfHl, line width=0.6pt, fill=bfHl, fill opacity=0.12},
        hlEdge/.style={draw=bfHl, line width=1.2pt},
        hlNonedge/.style={draw=bfHl, dash pattern=on 3pt off 2pt, line width=1.2pt},
    ]
    \def\bx{1.7}
    \def\rx{1.1}
    \def\ry{0.8}
    \def\off{0.7}
    \def\h{2.2}
    \def\panelsep{6.3}

    \pgfmathsetmacro{\bfunit}{\linewidth / (2*\panelsep + 2*(\bx+\rx))}
    \begin{scope}[x=\bfunit pt, y=\bfunit pt]

    \pgfmathsetmacro{\tphi}{atan2(\h/\ry, (2*\bx+\off)/\rx)}
    \pgfmathsetmacro{\tdelta}{acos(1/sqrt(((2*\bx+\off)/\rx)^2 + (\h/\ry)^2))}

    \def\bfSkipComplete{-1}
    \def\bfSkipAnti{-1}
    \def\bfSimLabel{$\sim$}

    \def\bfCrossingBase{%
        \coordinate (X) at (-\bx, 0);
        \coordinate (Y) at ( \bx, 0);
        \coordinate (x) at (-\bx-\off, 0.15);
        \coordinate (y) at ( \bx+\off, 0.15);
        \coordinate (x') at (-\bx-\off, \h);
        \coordinate (y') at ( \bx+\off, \h);

        \pgfmathsetmacro{\px}{2*\bx+\off}
        \pgfmathsetmacro{\fanfrom}{atan2(\ry*sin(\tphi+\tdelta) - \h, \rx*cos(\tphi+\tdelta) - \px)}
        \pgfmathsetmacro{\fanto}{atan2(\ry*sin(\tphi-\tdelta) - \h, \rx*cos(\tphi-\tdelta) - \px)}
        \foreach \i in {0,...,5} {
            \pgfmathsetmacro{\fanang}{\fanfrom + \i*(\fanto - \fanfrom)/5}
            \pgfmathsetmacro{\fanA}{(cos(\fanang))^2/\rx^2 + (sin(\fanang))^2/\ry^2}
            \pgfmathsetmacro{\fanB}{2*(\px*cos(\fanang)/\rx^2 + \h*sin(\fanang)/\ry^2)}
            \pgfmathsetmacro{\fanC}{\px^2/\rx^2 + \h^2/\ry^2 - 1}
            \pgfmathsetmacro{\fanlen}{(-(\fanB) - sqrt(max(0, \fanB*\fanB - 4*\fanA*\fanC)))/(2*\fanA)}
            \ifnum\i=\bfSkipComplete\relax\else
                \draw[complete] (y') -- ++(\fanang:\fanlen);
            \fi
            \ifnum\i=\bfSkipAnti\relax\else
                \draw[anticomplete, dash phase={mod(\i,2)*2.25pt}] (x') -- ++(180-\fanang:\fanlen);
            \fi
        }

        \draw[blob] (X) ellipse [x radius=\rx, y radius=\ry];
        \draw[blob] (Y) ellipse [x radius=\rx, y radius=\ry];
        \node[lbl, below=2pt] at ($(X) + (0, -\ry)$) {$X$};
        \node[lbl, below=2pt] at ($(Y) + (0, -\ry)$) {$Y$};

        \draw[sameType] (x') -- node[lbl, left=1pt, black] {\bfSimLabel} (x);
        \draw[sameType] (y') -- node[lbl, right=1pt, black] {\bfSimLabel} (y);

        \foreach \v in {x, y, x', y'}
            \node[vertex] at (\v) {};
        \node[lbl, below=2pt] at (x) {$x$};
        \node[lbl, below=2pt] at (y) {$y$};
        \node[lbl, above=2pt] at (x') {$x'$};
        \node[lbl, above=2pt] at (y') {$y'$};
    }

    \begin{scope}[name prefix=a-]
        \bfCrossingBase
    \end{scope}

    \begin{scope}[shift={(\panelsep, 0)}, name prefix=b-]
        \def\bfSkipAnti{3}
        \def\bfSimLabel{}
        \bfCrossingBase
        \pgfmathsetmacro{\bfsang}{180 - (\fanfrom + \bfSkipAnti*(\fanto - \fanfrom)/5)}
        \pgfmathsetmacro{\bfslen}{(2*\bx + \off - 0.4)/cos(\bfsang)}
        \coordinate (s) at ($(x') + (\bfsang:\bfslen)$);
        \draw[sampleSet] (s) ellipse [x radius=0.45, y radius=0.3];
        \draw[hlEdge] (s) -- (x);
        \draw[hlNonedge] (s) -- (x');
        \foreach \v in {x, x'}
            \node[vertex] at (\v) {};
        \node[vertex, fill=bfHl] at (s) {};
    \end{scope}

    \begin{scope}[shift={(2*\panelsep, 0)}, name prefix=c-]
        \def\bfSimLabel{}
        \bfCrossingBase
        \pgfmathsetmacro{\bfsang}{\fanfrom + 3*(\fanto - \fanfrom)/5}
        \pgfmathsetmacro{\bfslen}{-(2*\bx + \off - 0.4)/cos(\bfsang)}
        \coordinate (s) at ($(y') + (\bfsang:\bfslen)$);
        \draw[sampleSet] (s) ellipse [x radius=0.45, y radius=0.3];
        \draw[hlNonedge] (s) -- (y);
        \draw[hlEdge] (s) -- (y');
        \foreach \v in {y, y'}
            \node[vertex] at (\v) {};
        \node[vertex, fill=bfHl] at (s) {};
    \end{scope}

    \coordinate (bbtop) at (current bounding box.north);
    \coordinate (bbbot) at (current bounding box.south);
    \pgfresetboundingbox
    \path[use as bounding box]
        (-\bx-\rx,0 |- bbbot) rectangle (2*\panelsep+\bx+\rx,0 |- bbtop);
    \end{scope}
\end{tikzpicture}%
}}}%

        \caption{%
            Left: the forbidden configuration~\eqref{eq:ball-flip-crossing}.
            Middle: the case of the proof of \Cref{clm:dagger} where $\sample(X,Y)$ (marked in blue) lies in $Y$. 
            Right: the case where $\sample(X,Y)$ lies in $X$.}
        \label{fig:ball-flip-crossing}
    \end{center}

    Define the $S$-flip $H$ of $G$ where a pair $(P,Q)$ of parts of $\PP_S$ is flipped if and only if there is an edge in $G$ with endpoints in $P$ and in $Q$ and also in $A$ and in $B$. More formally, $H[P,Q]=G[P,Q]$ whenever
    $E(G) \cap PQ \cap AB = %\quad = \quad 
    \emptyset,$
    and otherwise, $H[P,Q]=\overline G[P,Q]$.
    

    The following three properties will yield the conclusion of the lemma:
    \begin{enumerate}[label=(P\arabic*), ref=P\arabic*]
        \item\label{itm:ball-flip-P1} $A$ and $\Ball_{r}^G(C) $ are nonadjacent in $H$.
        \hfill (proven in \Cref{clm:ball-flip-P1})
        \item\label{itm:ball-flip-P2} $\Ball_{r-1}^H(C)=\Ball_{r-1}^G(C)$.
        \hfill (proven in \Cref{clm:ball-flip-P2})
        \item\label{itm:ball-flip-P3}$\Ball_{r-1}^G(C)$ and $V(G) \setminus \Ball_{r}^G(C)$ are nonadjacent in $H$.
        \hfill (proven in \Cref{clm:ball-flip-P3})
    \end{enumerate}
    If $r = 0$, then already \ref{itm:ball-flip-P1} alone yields the conclusion. This justifies the implicit assumption $r \geq 1$ made in  \ref{itm:ball-flip-P2} and \ref{itm:ball-flip-P3}.
    By combining \ref{itm:ball-flip-P2} and \ref{itm:ball-flip-P3}, we derive:
    \begin{enumerate}[label=(P\arabic*), ref=P\arabic*,start = 4]
        \item\label{itm:ball-flip-P4}  $\Ball_{r}^H(C) \subseteq \Ball_{r}^G(C)$.
        % \hfill ()
    \end{enumerate}
    The conclusion of the lemma thus follows:
    \[
       A \cap \Ball_{r+1}^H(C)
       \  
    %    \underset{\rule{0pt}{8pt}\text{def.}}{=\vphantom{\subseteq}} 
    =
    \ 
       A \,\cap\, \Ball_1^H\big(\Ball_r^H(C)\big)
       \  \underset{\rule{0pt}{8pt}\text{\ref{itm:ball-flip-P4}}}{\subseteq} \ 
       A \cap \Ball_1^H\big(\Ball_r^G(C)\big)
       \  \underset{\rule{0pt}{8pt}\text{\ref{itm:ball-flip-P1}}}{=\vphantom{\subseteq}} \ 
       A \cap \Ball_r^G(C)
       \  \underset{\rule{0pt}{8pt}\text{ass.}}{=\vphantom{\subseteq}} \ 
       \emptyset.
    \]
    It only remains to establish \ref{itm:ball-flip-P1}--\ref{itm:ball-flip-P3}.
    


    \begin{claim}[\ref{itm:ball-flip-P1}]\label{clm:ball-flip-P1}
        $A$ and $B$ are nonadjacent in $H$.
    \end{claim}
    \begin{claimproof}
By contrapositive, suppose $ab\in E(H)$ for some $a \in A$ and $b \in B$.
    By definition of $H$,  $ab\in E(G)$ implies $ab\notin E(H)$, so
     $ab\notin E(G)$.
    Since $ab\in E(H)$, the parts of $\PP_S$ containing $a$ and $b$ are flipped, that is, there is some $a'b'\in E(G)$ with $a' \in A$ and $b' \in B$,
    and 
    \begin{center}
        $(a\sim a'$ and $b\sim b')$
        \qquad
        or
        \qquad
        $(a\sim b'$ and $b\sim a')$.
    \end{center}
    The first case cannot happen 
    as $\PP_S$ refines $\AA$ on $A$ and $\BB$ on $B$, but the two sets $\{a,a'\}$ and $\{b,b'\}$ are inhomogeneous as $ab\notin E(G)$ and $a'b'\in E(G)$.

    Thus, assume $a\sim b'$ and $b\sim a'$.
    Let $X,Y\in\AA$ contain $a',a$, respectively.
    Then $a'b'\in E(G)$ implies that $b'$ is complete to $X$, whereas  $ab\notin E(G)$ implies that $b$ is anticomplete to $Y$ (both by homogeneity). This contradicts \eqref{eq:ball-flip-crossing} with $(x,y,x',y')=(a',a,b,b')$.
\end{claimproof}

    For brevity, we define $D\coloneqq \Ball_{r-1}^G(C)$. Recall that $\PP_S(v)$ is the part of $\PP_S$ containing $v$.

    \begin{claim}\label{clm:disjoint-from-d}
        If $b$ is adjacent to some $a\in A$ in $G$ then $\PP_S(b)\cap D=\emptyset$.
    \end{claim}
    \begin{claimproof}
            The vertices in $B \cap \PP_S(b)$ have the same neighborhood on $A$ as~$b$. In particular, they are all adjacent to $a$, whereas $D \subseteq B$ is anticomplete to $A$ since $\Ball^G_1(D)=\Ball^G_r(C)$ is disjoint from $A$ by assumption.
    \end{claimproof}

    \begin{claim}[\ref{itm:ball-flip-P2}]\label{clm:ball-flip-P2}
        $\Ball_{r-1}^H(C)=D$.
    \end{claim}
    \begin{claimproof}
    The claim is trivially true for $r = 1$, hence assume $r\geq 2$.
    Consider any  $ab\in E(G)$ with $a \in A$ and $b \in B$ witnessing a flipped pair of $\PP_S$-parts $\PP_S(a)$ and $\PP_S(b)$.

    \begin{enumerate}
    \item \emph{$\PP_S(b)$ is disjoint from $\Ball_{r-2}^G(C)$:}
        
    This is implied by \Cref{clm:disjoint-from-d}, because $\Ball_{r-2}^G(C) \subseteq \Ball_{r-1}^G(C) =D$.

    \item \emph{$\PP_S(a)$ is disjoint from $\Ball_{r-2}^G(C)$:}

    Let $Q\in\BB$ contain $b$.
    Every vertex of $Q$ is adjacent to $a$ and hence lies in $B\setminus D$.
    Thus, $\rep(Q)\in S$ is adjacent to $a$ and nonadjacent to every vertex of $\Ball_{r-2}^G(C)$.
    Hence, $\rep(Q)$ witnesses that $a$ is in a different $\PP_S$-part than all the vertices in $\Ball_{r-2}^G(C)$. 
    \end{enumerate}

    This means that whenever two parts $\PP_S(a)$ and $\PP_S(b)$ are flipped, then neither of them intersects~$\Ball_{r-2}^G(C)$.
    Hence, the vertices in $\Ball_{r-2}^G(C)$ are incident to the same edges in $H$ as in $G$, which gives the desired conclusion 
    $\Ball_{r-1}^H(C) = \Ball_{r-1}^G(C)$.
    \end{claimproof}

    For brevity, we define $\overline B\coloneqq V(G)\setminus B$.

    \begin{claim}[\ref{itm:ball-flip-P3}]\label{clm:ball-flip-P3}
        $D$ and $\overline B$ are nonadjacent in $H$.
    \end{claim}

    \begin{claimproof}
            Otherwise, an edge $dz\in E(H)$ with $d \in D$ and $z \in \overline B$ is witnessed by some $ab\in E(G)$ with $a \in A$ and $b \in B$ (because $D$ and $\overline B$ are nonadjacent in $G$).
    By \Cref{clm:disjoint-from-d}, the $\PP_S$-part of $b$ is disjoint from $D$, so necessarily $d\sim a$ and $z\sim b$.
    Let $X\in\AA$ contain $a$.
    The vertex $b$ is complete to $X$, whereas $d$ is anticomplete to $\overline B$, contradicting \eqref{eq:ball-flip-crossing} for $(X,\overline B)$ with $(x,y,x',y')=(a,z,d,b)$.
    \end{claimproof}
    As described above, establishing \ref{itm:ball-flip-P1}--\ref{itm:ball-flip-P3} finishes the proof.
\end{proof}
% lax end


\subsection{Extension to tuples}\label{sec:closures-tuples}
In this subsection we prove \Cref{thm:closuresTuples}.
Recall that for $s\in\N$, an \emph{$s$-flip} of a graph $G$ is a $\PP$-flip of $G$ for some partition $\PP$ of $V(G)$ into at most $s$ parts.
For $S\subseteq V(G)$, the partition $\PP_S$ into atomic types over $S$ has at most $3^{|S|}$ parts, so every $s$-definable flip is a $3^s$-flip.
The flip metric $\dist_\PP$ associated with a partition $\PP$ of $V(G)$ is the maximum of $\dist^H$ over all $\PP$-flips $H$ of $G$, and $\Ball^\PP_r(v)\coloneqq\{u\in V(G):\dist_\PP(u,v)\leq r\}$ denotes its balls.
We use the same notation for the minimum distance between two tuples or sets.

% lax begin Lax185574.FlipMetricMonotone
% lax begin Lax185574Proofs.FlipMetricMonotone.isFlip_of_le
% lax begin Lax185574Proofs.FlipMetricMonotone.flipBall_subset_of_le
\begin{observation}\label{obs:partition-distance-monotone}
If $\QQ$ refines $\PP$, then every $\PP$-flip is also a $\QQ$-flip, and hence
\[
    \dist_\QQ(u,v)\geq \dist_\PP(u,v).
\]
\end{observation}
% lax end
% lax end
% lax end

We use the following metric-conversion result.    

% lax begin Lax185574.PartitionMetricConversion
% lax begin Lax185574Proofs.PartitionMetricConversion.exists_refinement_flip
\begin{lemma}[{\cite[Lemma 8]{bonnet2025separability-icalp}}]\label{lem:partition-metric-conversion}
    Let $G$ be a graph and let $\PP$ be a partition of $V(G)$. There is a refinement $\QQ$ of $\PP$ with $|\QQ|\leq |\PP|\cdot 2^{|\PP|}$ and a $\QQ$-flip $H$ of $G$ such that, for every $v\in V(G)$ and $r\in\N$,
    \[
        \Ball^H_r(v)\subseteq\Ball^\PP_{6r}(v).
    \]
\end{lemma}
% lax end
% lax end

We now combine the flips from  the Closure Theorem (\Cref{thm:closures}) for the individual entries of a tuple into a single flip.

% lax begin Lax185574.ClosureTheoremTuples
\thmClosuresTuples*
% lax end


% lax begin Lax185574Proofs.ClosureTheoremTuples.closure_theorem_tuples
\begin{proof}
    Let $c_0,s_0\in\N$ be given by \Cref{thm:closures} for $\CC$ and radius $6r$.
    Fix $G\in\CC$ and a finite set $A\subseteq V(G)$, and let $S$ and the sets $S(v)$ for $v\notin A$ be as in \Cref{thm:closures}.
    For a tuple $\bar v=(v_1,\ldots,v_k)\in(V(G)\setminus A)^k$, put
    \[
        S(\bar v)\coloneqq S(v_1)\cup\cdots\cup S(v_k).
    \]
    Then $S(\bar v)\subseteq S$ and $|S(\bar v)|\leq k\,s_0$.
    As $\PP_{S(\bar v)}$ refines every $\PP_{S(v_i)}$, \Cref{obs:partition-distance-monotone} gives
    \begin{equation}\label{eq:closures-tuples-far}
        \dist_{S(\bar v)}^G(v_i,A)\geq\dist_{S(v_i)}^G(v_i,A)>6r
        \qquad\text{for every }i\in[k].
    \end{equation}

    For every $T\subseteq S$ with $|T|\leq ks_0$, the partition $\PP_T$ has at most $3^{ks_0}$ parts.
    So \Cref{lem:partition-metric-conversion} yields an $s$-flip $H_T$ of $G$, for some $s\leq O_{\CC,r,k}(1)$ independent of $T$, such that
    \[
        \Ball^{H_T}_r(x)\subseteq\Ball^{T}_{6r}(x)
        \qquad\text{for every }x\in V(G).
    \]
    Let $\HH\coloneqq\{H_T \,:\, T\subseteq S,\ |T|\leq ks_0\}$.
    Then $|\HH|\leq(|S|+1)^{ks_0}\leq O_{\CC,r,k}(|A|^{c})$ for $c\coloneqq c_0\, k\, s_0$.

    Finally, let $\bar v=(v_1,\ldots,v_k)\in(V(G)\setminus A)^k$ and $H\coloneqq H_{S(\bar v)}\in\HH$.
    For every $i\in[k]$, the ball $\Ball^H_r(v_i)\subseteq\Ball^{S(\bar v)}_{6r}(v_i)$ is disjoint from $A$ by \eqref{eq:closures-tuples-far}.
    Thus $\dist^H(\bar v,A)>r$.
\end{proof}
% lax end


\section{Number of types over a set of parameters}\label{sec:types-over-sets}
In this section we prove \Cref{thm:type-bound-dp}, which bounds the number of types of $k$-tuples over a set $U$ of parameters polynomially in $|U|$. This also yields \Cref{thm:type-bound}, by \Cref{lem:mondep-dpmin}.

We use a consequence of locality of first-order logic, which bounds the number of types of tuples that are far from $U$. 
% Afterwards, we reduce the general case to this one using the Closure Theorem for Tuples.
% \subsection{Locality of first-order logic}
% \label{sec:locality-of-fo}
% \newcommand{\adj}{\mathrm{adj}}
% In this subsection we show that, in a graph with boundedly many unary predicates, the tuples far from a set $U$ realize only boundedly many types over $U$, with a bound independent of $|U|$ (\Cref{lem:sep-types,lem:sep-types-tuple}).
The following result is an adaptation of~\cite[Corollary~16]{number-of-types}. We provide a self-contained proof in \Cref{app:gaifman}.



% lax begin Lax185574.SeparatingTypesTuples
\begin{restatable}{lemma}{septypestuples}\label{lem:sep-types-tuple}
    For every $q,k,c\in\N$ there are numbers $r=r(q)$ and $T=T(q,k,c)$ with the following property. Let $H$ be a graph expanded by at most $c$ unary predicates, and let $U\subseteq V(H)$ be finite. Fix a tuple $\bar u$ of vertices of $U$ of finite length.
    Let $F \coloneqq \{a\in V(H) \,:\, \dist^H(a,U) > r\}$. 
    Then
    \[
      \left|
        \{
        \tp_q^H(\bar a\bar u/U) \,:\, \bar a \in F^k  
        \}
      \right|
      \leq
      T(q,k,c).      
    \]
\end{restatable}
% lax end









% \subsection{Type bounds}

We now prove \Cref{thm:type-bound-dp} by combining \Cref{lem:sep-types-tuple} with the Closure Theorem for Tuples (\Cref{thm:closuresTuples}).



% lax begin Lax185574.DpMinimalTypeBound
{\typeboundgeneraltrue\dpmintypes*}
% lax end

% lax begin Lax185574Proofs.DpMinimalTypeBound.typeCountOver_le
\begin{proof}
    Fix $G\in\CC$, its expansion $G^+$ by $c$ unary predicates, and a finite set $U\subseteq V(G)$.
    Let $r=r(q)$ and $T$ be given by \Cref{lem:sep-types-tuple}, and put $\overline{U}\coloneqq V(G)\setminus U$.
    By the Closure Theorem for Tuples (\Cref{thm:closuresTuples}) for $\CC$, radius $r$, and $k$, there is some $s\leq O_{\CC,q,k}(1)$ and some set $\HH$ of $\poly_{\CC,q,k}(|U|)$ many $s$-flips of $G$ such that every $\bar v\in\overline{U}^k$ has some $H(\bar v)\in\HH$ with $\dist^{H(\bar v)}(\bar v,U)>r$.
    Note that $s$ and $\HH$ do not depend on the unary predicates of $G^+$.
    For each $H\in\HH$, fix a partition $\QQ_H$ of $V(G)$ into at most $s$ parts such that $H$ is a $\QQ_H$-flip of $G$, and let $H^+$ be the expansion of $H$ by the $c$ unary predicates of $G^+$ and by at most $s$ unary predicates marking the parts of $\QQ_H$.
    The edge relation of $G$ is quantifier-free definable in $H^+$, by flipping back the pairs of parts that are flipped in $H$.
    Substituting this definition for every edge atom turns each formula over $G^+$ with parameters from $U$ into a formula over $H^+$ with the same parameters and quantifier rank that is satisfied by the same tuples.
    Hence, for every $H$ and tuple $\bar w$ of vertices,
    \begin{center}
        $\tp_q^{G^+}(\bar w/U)$ depends only on $\tp_q^{H^+}(\bar w/U)$.
    \end{center}

    We want to bound the number of types of tuples $\bar v \in V(G)^k$ over $U$.
    By padding and permuting $2k$-tuples, up to a multiplicative factor of $O_k(1)$, it suffices to bound the number of types for the slightly longer tuples of the form $\bar a\bar u$ with $\bar a \in \overline{U}^k$ and $\bar u \in U^k$.
    
    For $H\in\HH$, let $F_H\coloneqq\{w\in V(G) \,:\, \dist^H(w,U)>r\} \subseteq \overline{U}$.
    As $\bar a\in\overline{U}^k$ lies in $F_{H(\bar a)}^k$ for some $H(\bar a) \in \HH$, we have
    \[
        \{\tp_q^{G^+}(\bar a\bar u/U) \,:\, \bar a\in\overline{U}^k,\ \bar u\in U^k\}
        =
        \bigcup_{H\in\HH}\ \bigcup_{\bar u\in U^k}
        \{\tp_q^{G^+}(\bar a\bar u/U) \,:\, \bar a\in F_H^k\}.
    \]
    It remains to bound each of the $|\HH|\cdot|U|^k$ terms on the right.
    So fix $H\in\HH$ and $\bar u\in U^k$.
    As $H^+$ has at most $c+s$ unary predicates, \Cref{lem:sep-types-tuple}, applied to $H^+$, $U$, and $\bar u$, gives
    \[
        \left|\{\tp_q^{H^+}(\bar a\bar u/U) \,:\, \bar a\in F_H^k\}\right|
        \leq T(q,k,c+s).
    \]
    Since $\tp_q^{G^+}(\bar a\bar u/U)$ depends only on $\tp_q^{H^+}(\bar a\bar u/U)$, the same bound holds with $G^+$ in place of $H^+$.
    Hence,
    \[
        \left|\{\tp_q^{G^+}(\bar a\bar u/U) \,:\, \bar a\in\overline{U}^k,\ \bar u\in U^k\}\right|
        \leq|\HH|\cdot|U|^k\cdot T(q,k,c+s)
        \leq O_{\CC,q,k,c}(1)\cdot\poly_{\CC,q,k}(|U|),
    \]
    where we use $T(q,k,c+s)\leq O_{\CC,q,k,c}(1)$, as $s\leq O_{\CC,q,k}(1)$, and $|\HH|\leq\poly_{\CC,q,k}(|U|)$.
\end{proof}
% lax end


